\chapter{Backend Implementation}

\vspace{-0.5cm}
\begin{center}
\begin{minipage}{0.85\textwidth}
\small
\textbf{Contents}\\[0.3cm]
\hspace*{1cm} 4.1 \quad Overview \dotfill \pageref{sec:ch4overview}\\
\hspace*{1cm} 4.2 \quad Objectives \dotfill \pageref{sec:ch4objectives}\\
\hspace*{1cm} 4.3 \quad Scope \dotfill \pageref{sec:ch4scope}\\
\hspace*{1cm} 4.4 \quad System Architecture \dotfill \pageref{sec:ch4arch}\\
\hspace*{1cm} 4.5 \quad Database Design \dotfill \pageref{sec:ch4db}\\
\hspace*{1cm} 4.6 \quad API Documentation \dotfill \pageref{sec:ch4api}\\
\hspace*{1cm} 4.7 \quad AI Microservices \dotfill \pageref{sec:ch4ai}\\
\hspace*{1cm} 4.8 \quad Admin Panel (FilamentPHP 3) \dotfill \pageref{sec:ch4admin}\\
\hspace*{1cm} 4.9 \quad Testing \dotfill \pageref{sec:ch4testing}\\
\hspace*{1cm} 4.10 \quad Challenges Faced \dotfill \pageref{sec:ch4challenges}\\
\hspace*{1cm} 4.11 \quad Future Work \dotfill \pageref{sec:ch4future}\\
\end{minipage}
\end{center}
\vspace{0.5cm}
\newpage

% ===========================
\section{Overview}
\label{sec:ch4overview}

The modern reader, student, and researcher face a growing challenge: the sheer
volume of available books and digital resources makes it increasingly difficult
to discover content that truly matches one's interests, academic needs, or
research goals. Traditional library systems rely on keyword-based search, which
often returns irrelevant results and fails to understand the user's intent.
There is also a lack of platforms that combine book management, educational
courses, intelligent recommendations, and AI-powered features in one unified
system.

Learnova addresses these challenges by building a centralized, scalable, and
AI-integrated backend that powers personalized book and course discovery,
semantic search, book summarization, and direct Q\&A from book content ---
all through a clean RESTful API consumed by the Flutter mobile application.

% ===========================
\section{Objectives}
\label{sec:ch4objectives}

The main objectives of the backend system are:

\begin{itemize}
  \item To provide a secure RESTful API consumed by the Flutter mobile
        application.
  \item To implement full CRUD operations for books, courses, authors,
        categories, and users.
  \item To integrate JWT-based authentication and role-based access control
        (Admin, Researcher, Student).
  \item To support a review and rating system allowing users to evaluate books.
  \item To manage user bookmark collections, personal reading shelves, and
        favorites.
  \item To integrate 6 Python AI microservices covering ingestion,
        recommendations, semantic search, summarization, and Q\&A.
  \item To support course management with video lessons and AI-powered lesson
        transcription.
  \item To deliver an admin panel (FilamentPHP 3) for efficient management of
        all platform resources.
  \item To implement OTP-based password reset and social login
        (Google/Facebook).
\end{itemize}

% ===========================
\section{Scope}
\label{sec:ch4scope}

\subsection{Included}

\begin{itemize}
  \item User registration, login, OTP password reset, and social login
        (Google/Facebook).
  \item JWT-based authentication via Laravel Sanctum with role-based access
        control (admin, researcher, student).
  \item Full CRUD for Books, Courses, Course Videos, Authors, Categories, and
        Chapters.
  \item Advanced search and filtering: global search, semantic search, search
        history, and trending.
  \item Review and rating system for books.
  \item Bookmark, favorites, and personal library shelves management.
  \item Onboarding: user topic preferences and profile setup.
  \item Admin panel (FilamentPHP 3) with full CRUD, file uploads, and analytics
        dashboard.
  \item Full AI integration: 6 Python microservices via Docker (ingestion,
        embedding, recommendations, semantic search, summarization, RAG Q\&A).
  \item AI-powered lesson transcription via \texttt{TranscribeAllLessons}
        command.
  \item Proxy architecture: Laravel acts as a gateway forwarding requests to
        AI microservices.
\end{itemize}

\subsection{Excluded}

\begin{itemize}
  \item Video conferencing or live streaming functionality.
  \item Payment processing and subscription management.
  \item Direct instructor consultation chat.
  \item OCR for scanned documents (planned future enhancement).
  \item Audio-book streaming (planned future enhancement).
\end{itemize}

% ===========================
\section{System Architecture}
\label{sec:ch4arch}

The Learnova backend follows a layered RESTful API architecture built on
\textbf{Laravel 12}. The system is designed around a proxy pattern where
Laravel serves as the central gateway: it handles authentication, business
logic, and database operations, while forwarding AI-related requests to
dedicated Python microservices running in Docker containers.

\begin{table}[H]
\centering
\caption{System Architecture Layers and Responsibilities}
\label{tab:arch}
\begin{tabularx}{\textwidth}{|l|l|X|}
\hline
\textbf{Layer} & \textbf{Technology} & \textbf{Responsibility} \\
\hline
Mobile Frontend    & Flutter (Dart)         & Cross-platform UI, state management, API consumption \\
\hline
Backend API        & Laravel 12 (PHP 8.1+)  & RESTful endpoints, business logic, authentication, AI proxy \\
\hline
Admin Panel        & FilamentPHP 3          & Admin UI, CRUD management, analytics dashboard \\
\hline
Database           & MySQL 8.0              & Persistent storage for users, books, courses, reviews \\
\hline
AI --- Ingestion   & Python (Port 8000)     & PDF/EPUB upload, text extraction, chunking, Kafka \\
\hline
AI --- Embedding   & Python (Background)    & Kafka consumer, vector generation, Pinecone storage \\
\hline
AI --- Recommendations & Python (Port 8001) & NCF + CBF hybrid recommendation engine \\
\hline
AI --- Semantic Search & Python (Port 8002) & Semantic search and result re-ranking \\
\hline
AI --- Summarization   & Python (Port 8003) & FLAN-T5 book summarization \\
\hline
AI --- RAG Q\&A    & Python (Port 8004)     & Groq LLM Q\&A with citations (sync + SSE streaming) \\
\hline
\end{tabularx}
\end{table}

\subsection{AI Proxy Architecture}

All AI-related API calls follow a proxy pattern. The mobile app sends requests
to Laravel, which forwards them to the appropriate AI microservice and returns
the result. If an AI service is unavailable, Laravel returns a graceful error
message rather than crashing the application. Figure~\ref{fig:proxy} illustrates
this flow.

\begin{figure}[H]
\centering
\begin{tabular}{ccc}
\textbf{Mobile App} & \textbf{Laravel API} & \textbf{AI Microservices} \\[6pt]
\multicolumn{3}{l}{POST /api/ai/search $\longrightarrow$ POST /search $\longrightarrow$ search-service:8002} \\
\multicolumn{3}{l}{POST /api/ai/recommend $\longrightarrow$ POST /recommendations $\longrightarrow$ rec-service:8001} \\
\multicolumn{3}{l}{POST /api/ai/summarise $\longrightarrow$ POST /summarise $\longrightarrow$ summarise-service:8003} \\
\multicolumn{3}{l}{POST /api/ai/qa/stream $\longrightarrow$ POST /qa/stream $\longrightarrow$ rag-service:8004} \\
\end{tabular}
\caption{AI Proxy Request Flow in Learnova}
\label{fig:proxy}
\end{figure}

% ===========================
\section{Database Design}
\label{sec:ch4db}

The system uses \textbf{MySQL 8.0} as the primary relational database. The
schema is normalized with clear foreign-key relationships between all entities.
Laravel 12 migrations are used to version-control the schema. The database
covers user management, book and course content, AI metadata, community
features, and personal library management.

\subsection{users}

\begin{table}[H]
\centering
\caption{users Table Schema}
\begin{tabularx}{\textwidth}{|l|l|X|}
\hline
\textbf{Column} & \textbf{Type} & \textbf{Description} \\
\hline
id                   & BIGINT (PK)    & Auto-incremented primary key \\
\hline
name                 & VARCHAR(255)   & Full name of the user \\
\hline
email                & VARCHAR(255)   & Unique email address \\
\hline
password             & VARCHAR(255)   & Bcrypt-hashed password \\
\hline
role                 & ENUM           & Role: admin, researcher, student \\
\hline
email\_verified\_at  & TIMESTAMP      & Email verification timestamp \\
\hline
created\_at / updated\_at & TIMESTAMP & Record timestamps \\
\hline
\end{tabularx}
\end{table}

\subsection{books}

\begin{table}[H]
\centering
\caption{books Table Schema}
\begin{tabularx}{\textwidth}{|l|l|X|}
\hline
\textbf{Column} & \textbf{Type} & \textbf{Description} \\
\hline
id             & BIGINT (PK)   & Auto-incremented primary key \\
\hline
title          & VARCHAR(255)  & Book title \\
\hline
isbn           & VARCHAR(255)  & Unique ISBN (used as AI service key) \\
\hline
author\_id     & BIGINT (FK)   & Reference to authors table \\
\hline
category\_id   & BIGINT (FK)   & Reference to categories table \\
\hline
description    & TEXT          & Book synopsis \\
\hline
cover\_image   & VARCHAR(255)  & Path to uploaded cover image \\
\hline
file\_path     & VARCHAR(255)  & Path to PDF/EPUB file for AI ingestion \\
\hline
published\_year & YEAR         & Year of publication \\
\hline
is\_free       & BOOLEAN       & Whether the book is free or paid \\
\hline
views\_count   & INT           & View counter for trending analytics \\
\hline
created\_at / updated\_at & TIMESTAMP & Record timestamps \\
\hline
\end{tabularx}
\end{table}

\subsection{courses}

\begin{table}[H]
\centering
\caption{courses Table Schema}
\begin{tabularx}{\textwidth}{|l|l|X|}
\hline
\textbf{Column} & \textbf{Type} & \textbf{Description} \\
\hline
id             & BIGINT (PK)   & Auto-incremented primary key \\
\hline
title          & VARCHAR(255)  & Course title \\
\hline
description    & TEXT          & Course description \\
\hline
category\_id   & BIGINT (FK)   & Reference to categories table \\
\hline
cover\_image   & VARCHAR(255)  & Course thumbnail image \\
\hline
is\_free       & BOOLEAN       & Whether the course is free or paid \\
\hline
is\_featured   & BOOLEAN       & Flag for featured courses \\
\hline
created\_at / updated\_at & TIMESTAMP & Record timestamps \\
\hline
\end{tabularx}
\end{table}

\subsection{course\_videos}

\begin{table}[H]
\centering
\caption{course\_videos Table Schema}
\begin{tabularx}{\textwidth}{|l|l|X|}
\hline
\textbf{Column} & \textbf{Type} & \textbf{Description} \\
\hline
id          & BIGINT (PK)   & Auto-incremented primary key \\
\hline
course\_id  & BIGINT (FK)   & Reference to courses table \\
\hline
title       & VARCHAR(255)  & Video lesson title \\
\hline
video\_url  & VARCHAR(255)  & URL or path to video file \\
\hline
transcript  & TEXT          & AI-generated transcript (TranscribeAllLessons) \\
\hline
duration    & INT           & Video duration in seconds \\
\hline
order       & INT           & Display order within the course \\
\hline
created\_at / updated\_at & TIMESTAMP & Record timestamps \\
\hline
\end{tabularx}
\end{table}

\subsection{reviews}

\begin{table}[H]
\centering
\caption{reviews Table Schema}
\begin{tabularx}{\textwidth}{|l|l|X|}
\hline
\textbf{Column} & \textbf{Type} & \textbf{Description} \\
\hline
id        & BIGINT (PK)     & Auto-incremented primary key \\
\hline
user\_id  & BIGINT (FK)     & Reference to users table \\
\hline
book\_id  & BIGINT (FK)     & Reference to books table \\
\hline
rating    & TINYINT (1--5)  & Star rating \\
\hline
comment   & TEXT            & Optional review text \\
\hline
created\_at / updated\_at & TIMESTAMP & Record timestamps \\
\hline
\end{tabularx}
\end{table}

\subsection{bookmarks and personal\_shelves}

\begin{table}[H]
\centering
\caption{bookmarks / personal\_shelves Table Schema}
\begin{tabularx}{\textwidth}{|l|l|X|}
\hline
\textbf{Column} & \textbf{Type} & \textbf{Description} \\
\hline
id           & BIGINT (PK)  & Auto-incremented primary key \\
\hline
user\_id     & BIGINT (FK)  & Reference to users table \\
\hline
book\_id     & BIGINT (FK)  & Reference to books table \\
\hline
shelf\_type  & ENUM         & Type: reading, completed, favorite \\
\hline
created\_at / updated\_at & TIMESTAMP & Record timestamps \\
\hline
\end{tabularx}
\end{table}

\subsection{ai\_metadata}

\begin{table}[H]
\centering
\caption{ai\_metadata Table Schema}
\begin{tabularx}{\textwidth}{|l|l|X|}
\hline
\textbf{Column} & \textbf{Type} & \textbf{Description} \\
\hline
id              & BIGINT (PK)  & Auto-incremented primary key \\
\hline
book\_id        & BIGINT (FK)  & Reference to books table \\
\hline
summary         & TEXT         & AI-generated summary (FLAN-T5) \\
\hline
keywords        & JSON         & Extracted keywords array \\
\hline
vector\_status  & ENUM         & Embedding status: pending, processing, done \\
\hline
mindmap\_data   & JSON         & Optional visual mind map data \\
\hline
created\_at / updated\_at & TIMESTAMP & Record timestamps \\
\hline
\end{tabularx}
\end{table}

% ===========================
\section{API Documentation}
\label{sec:ch4api}

The backend exposes a comprehensive RESTful API. All endpoints return JSON
responses. Authentication is handled via Laravel Sanctum Bearer tokens passed
in the \texttt{Authorization} header. The base URL for all API calls is
\texttt{http://localhost:8000/api}.

\subsection{Authentication Endpoints}

\begin{table}[H]
\centering
\caption{Authentication API Endpoints}
\begin{tabularx}{\textwidth}{|l|l|X|}
\hline
\textbf{Method} & \textbf{Endpoint} & \textbf{Description} \\
\hline
POST & /api/register                    & Register a new user account \\
\hline
POST & /api/login                       & Login and receive Sanctum token \\
\hline
POST & /api/logout                      & Invalidate current token \\
\hline
POST & /api/forgot-password/send-otp    & Send OTP for password reset \\
\hline
POST & /api/forgot-password/verify-otp  & Verify the received OTP code \\
\hline
POST & /api/forgot-password/reset       & Reset password after OTP verification \\
\hline
GET  & /api/auth/\{provider\}/redirect  & Redirect to Google or Facebook OAuth \\
\hline
GET  & /api/auth/\{provider\}/callback  & OAuth callback handler (social login) \\
\hline
\end{tabularx}
\end{table}

\subsection{Books and Discovery Endpoints}

\begin{table}[H]
\centering
\caption{Books and Discovery API Endpoints}
\begin{tabularx}{\textwidth}{|l|l|X|}
\hline
\textbf{Method} & \textbf{Endpoint} & \textbf{Description} \\
\hline
GET    & /api/home                    & Home page: personalized recommendations \\
\hline
GET    & /api/discover                & Discover page: trending and featured content \\
\hline
GET    & /api/discover/trending-20    & Top 20 trending books by views count \\
\hline
GET    & /api/discover/free-20        & Top 20 free books \\
\hline
GET    & /api/search                  & General search across books and content \\
\hline
GET    & /api/search/global           & Global comprehensive search \\
\hline
GET    & /api/search/filter/\{type\}  & Filtered search by type (book/author/course) \\
\hline
DELETE & /api/search/history/\{id?\}  & Delete search history entry \\
\hline
POST   & /api/books/\{id\}/favorite   & Toggle book favorite status \\
\hline
POST   & /api/books/\{id\}/listen     & Record a book listen/view event \\
\hline
GET    & /api/books/\{id\}/share      & Get shareable link for a book \\
\hline
\end{tabularx}
\end{table}

\subsection{Courses Endpoints}

\begin{table}[H]
\centering
\caption{Courses API Endpoints}
\begin{tabularx}{\textwidth}{|l|l|X|}
\hline
\textbf{Method} & \textbf{Endpoint} & \textbf{Description} \\
\hline
GET    & /api/courses                            & List all courses with pagination \\
\hline
GET    & /api/courses/featured                   & Get featured courses \\
\hline
GET    & /api/courses/free                       & Get free courses \\
\hline
GET    & /api/courses/category/\{id\}            & Filter courses by category \\
\hline
GET    & /api/courses/\{id\}                     & Get course details including video list \\
\hline
POST   & /api/courses                            & Create a new course (Admin only) \\
\hline
PUT    & /api/courses/\{id\}                     & Update course details (Admin only) \\
\hline
DELETE & /api/courses/\{id\}                     & Delete a course (Admin only) \\
\hline
GET    & /api/courses/\{id\}/videos              & Get all videos for a course \\
\hline
POST   & /api/courses/\{id\}/videos              & Add a new video to a course \\
\hline
PUT    & /api/courses/\{id\}/videos/\{vidId\}    & Update a course video \\
\hline
DELETE & /api/courses/\{id\}/videos/\{vidId\}    & Delete a course video \\
\hline
\end{tabularx}
\end{table}

\subsection{AI Services Endpoints}

All AI endpoints follow a proxy pattern: Laravel receives the request, validates
authentication, and forwards to the appropriate Python microservice. If a
service is unavailable, a graceful error is returned.

\begin{table}[H]
\centering
\caption{AI Services API Endpoints}
\begin{tabularx}{\textwidth}{|l|l|l|X|}
\hline
\textbf{Method} & \textbf{Endpoint} & \textbf{Service} & \textbf{Description} \\
\hline
GET  & /api/ai/health              & All       & Health check for all AI microservices \\
\hline
POST & /api/ai/ingest              & :8000     & Upload PDF/EPUB for processing \\
\hline
GET  & /api/ai/ingest/\{isbn\}/status & :8000  & Track ingestion pipeline progress \\
\hline
POST & /api/ai/recommendations     & :8001     & Get personalized book recommendations \\
\hline
GET  & /api/ai/recommendations/cold-start & :8001 & Cold-start recommendations \\
\hline
POST & /api/ai/search/semantic     & :8002     & Semantic search by query \\
\hline
POST & /api/ai/search/similar/\{isbn\} & :8002 & Find similar books by ISBN \\
\hline
POST & /api/ai/summarise           & :8003     & Summarize a book by ISBN \\
\hline
GET  & /api/ai/summarise/\{isbn\}/progress & :8003 & Track summarization progress \\
\hline
POST & /api/ai/qa/sync             & :8004     & Full Q\&A response with citations \\
\hline
POST & /api/ai/qa/stream           & :8004     & Streaming Q\&A via SSE tokens \\
\hline
\end{tabularx}
\end{table}

% ===========================
\section{AI Microservices}
\label{sec:ch4ai}

The AI layer consists of 6 Python microservices deployed via Docker Compose.
Each service is independently scalable and communicates through well-defined
REST interfaces. The services use \textbf{Kafka} for asynchronous pipeline
processing and \textbf{Pinecone} as the vector database for semantic search
capabilities.

\begin{table}[H]
\centering
\caption{Python AI Microservices Overview}
\begin{tabularx}{\textwidth}{|l|l|l|X|}
\hline
\textbf{Service} & \textbf{Port} & \textbf{Model/Tech} & \textbf{Function} \\
\hline
ingestion        & 8000 & PyMuPDF + Kafka          & PDF/EPUB upload, text extraction, chunking \\
\hline
embedding        & ---  & Kafka + Pinecone         & Background: generate vectors, store in Pinecone \\
\hline
rec\_service     & 8001 & NCF + CBF hybrid         & Personalized recommendations \\
\hline
search\_service  & 8002 & Hugging Face + re-ranker & Semantic search with intent understanding \\
\hline
summarise\_service & 8003 & FLAN-T5              & Multi-type summarization: short, detailed, mind map \\
\hline
rag\_service     & 8004 & Groq LLM               & Q\&A from book text with citations (sync + SSE) \\
\hline
\end{tabularx}
\end{table}

\subsection{Courses AI --- TranscribeAllLessons}

As part of the courses AI integration, a dedicated Laravel Artisan command
\texttt{TranscribeAllLessons} was developed. This command processes course video
lessons by sending them to the AI transcription pipeline, generating text
transcripts stored in the \texttt{course\_videos} table. This enables
searchable lesson content, AI-powered lesson summaries, and accessibility
improvements for hearing-impaired users.

\begin{lstlisting}[language=bash, caption=Running the TranscribeAllLessons command]
# Run transcription for all unprocessed lessons
php artisan transcribe:all-lessons

# The command:
# 1. Fetches all course_videos without transcripts
# 2. Sends each video to the AI transcription service
# 3. Stores the returned transcript in the database
# 4. Logs progress and handles failures gracefully
\end{lstlisting}

\subsection{AI Service Configuration}

All AI service URLs and timeout settings are centralized in
\texttt{config/ai.php}. The \texttt{AiService.php} class in
\texttt{app/Services} provides a unified interface for all AI microservice
calls, handling connection errors gracefully.

\begin{table}[H]
\centering
\caption{AI Service Configuration Parameters}
\begin{tabularx}{\textwidth}{|l|l|X|}
\hline
\textbf{Config Key} & \textbf{Default Value} & \textbf{Description} \\
\hline
ai.ingestion\_url & http://localhost:8000 & Ingestion service base URL \\
\hline
ai.rec\_url       & http://localhost:8001 & Recommendation service base URL \\
\hline
ai.search\_url    & http://localhost:8002 & Semantic search service base URL \\
\hline
ai.summarise\_url & http://localhost:8003 & Summarization service base URL \\
\hline
ai.rag\_url       & http://localhost:8004 & RAG Q\&A service base URL \\
\hline
ai.timeout        & 30 seconds            & Default HTTP timeout for AI requests \\
\hline
\end{tabularx}
\end{table}

% ===========================
\section{Admin Panel (FilamentPHP 3)}
\label{sec:ch4admin}

The backend includes a full-featured admin panel built with \textbf{FilamentPHP
3}, accessible at \texttt{/admin}. The panel is organized into four sections
covering the entire platform.

\begin{table}[H]
\centering
\caption{Admin Panel Sections and Resources}
\begin{tabularx}{\textwidth}{|l|l|X|}
\hline
\textbf{Section} & \textbf{Resources} & \textbf{Description} \\
\hline
Library        & Books, Courses, Authors, Categories, Chapters & Full CRUD with file uploads \\
\hline
Community      & Reviews, Personal Shelves & User-generated content moderation \\
\hline
AI Engine      & AI Metadata & View summaries, keywords, vector status \\
\hline
Administration & Users       & User management with role assignment \\
\hline
\end{tabularx}
\end{table}

\begin{table}[H]
\centering
\caption{Admin Panel Default Credentials}
\begin{tabularx}{\textwidth}{|l|X|}
\hline
\textbf{Field} & \textbf{Value} \\
\hline
Admin Panel URL   & http://127.0.0.1:8000/admin \\
\hline
Default Email     & admin@admin.com \\
\hline
Default Password  & password \\
\hline
Note & Change credentials immediately after first deployment \\
\hline
\end{tabularx}
\end{table}

% ===========================
\section{Testing}
\label{sec:ch4testing}

The backend APIs were tested using \textbf{Postman} for manual endpoint
validation and Laravel's built-in \textbf{PHPUnit} framework for automated
tests. Each endpoint was verified for correct HTTP response codes, data
accuracy, authentication enforcement, and error handling for invalid inputs.

\subsection{Testing Tools}

\begin{table}[H]
\centering
\caption{Testing Tools and Scope}
\begin{tabularx}{\textwidth}{|l|X|}
\hline
\textbf{Tool} & \textbf{Scope} \\
\hline
Postman           & Manual API testing: auth flows, CRUD operations, AI endpoints, edge cases \\
\hline
PHPUnit           & Automated unit and integration tests for models and API responses \\
\hline
Browser (Manual)  & Admin panel CRUD, file uploads, dashboard analytics \\
\hline
Docker Health Checks & Verify all 6 AI microservices via \texttt{/health} endpoints \\
\hline
\end{tabularx}
\end{table}

\subsection{Sample Test Cases}

\begin{table}[H]
\centering
\caption{Sample API Test Cases and Expected Results}
\begin{tabularx}{\textwidth}{|l|l|l|X|}
\hline
\textbf{Endpoint} & \textbf{Method} & \textbf{Test Case} & \textbf{Expected Result} \\
\hline
/api/register        & POST & Valid data            & 201 Created + User object \\
\hline
/api/register        & POST & Duplicate email       & 422 Unprocessable Entity \\
\hline
/api/login           & POST & Valid credentials     & 200 OK + Sanctum token \\
\hline
/api/login           & POST & Wrong password        & 401 Unauthorized \\
\hline
/api/forgot-password/send-otp & POST & Valid phone & 200 OK + OTP sent \\
\hline
/api/discover/trending-20 & GET & Fetch trending   & 200 OK + 20 books \\
\hline
/api/courses         & GET  & List all courses      & 200 OK + Paginated \\
\hline
/api/ai/health       & GET  & All services running  & 200 OK + All healthy \\
\hline
/api/ai/search/semantic & POST & Valid query        & 200 OK + Ranked results \\
\hline
/api/ai/recommendations & POST & Valid user\_id    & 200 OK + Recommended books \\
\hline
/api/ai/summarise    & POST & Valid ISBN + type     & 200 OK + Summary text \\
\hline
/api/ai/qa/sync      & POST & Question about book   & 200 OK + Answer + citations \\
\hline
/api/bookmarks       & GET  & No auth token         & 401 Unauthorized \\
\hline
\end{tabularx}
\end{table}

% ===========================
\section{Challenges Faced}
\label{sec:ch4challenges}

\begin{itemize}
  \item \textbf{AI Microservices Integration}: Designing a proxy architecture
        in Laravel that gracefully handles failures from any of the 6 Python
        microservices without crashing the main application required careful
        error handling and timeout configuration in \texttt{AiService.php}.

  \item \textbf{Docker Environment Configuration}: Setting up Kafka, Pinecone,
        and 6 containerized services with correct environment variables and
        ensuring inter-container networking required extensive troubleshooting
        of the \texttt{docker-compose.yml} configuration.

  \item \textbf{Courses AI --- TranscribeAllLessons}: Developing the Artisan
        command to batch-process course video lessons through the AI
        transcription pipeline required handling partial failures, retry logic,
        and efficient progress logging for large course libraries.

  \item \textbf{Dual Authentication System}: Configuring Laravel Sanctum for
        API JWT tokens alongside FilamentPHP's session-based admin
        authentication without conflicts required careful guard configuration
        in \texttt{auth.php}.

  \item \textbf{Social Login Integration}: Implementing Google and Facebook
        OAuth flows with proper callback handling, user creation/linking logic,
        and token generation for the mobile app required careful configuration
        of Laravel Socialite.

  \item \textbf{Role-Based Access Control}: Implementing middleware that
        correctly restricts admin-only endpoints while keeping public discovery
        endpoints open required thorough testing across three user roles.

  \item \textbf{Database Migration Ordering}: Managing foreign-key dependencies
        across multiple related tables required careful ordering of Laravel
        migrations to avoid constraint violations.
\end{itemize}

% ===========================
\section{Future Work}
\label{sec:ch4future}

\begin{itemize}
  \item Merge the courses AI branch (\texttt{TranscribeAllLessons}) into main
        and extend to support real-time transcription during video upload.
  \item Add audio-book streaming with AI-generated narration using
        text-to-speech models.
  \item Implement OCR for scanned document digitization using Tesseract or
        similar tools.
  \item Develop an advanced admin analytics dashboard with user engagement
        metrics, book popularity trends, and AI usage statistics.
  \item Expand the recommendation engine to include course recommendations
        alongside book recommendations.
  \item Add multi-language support (Arabic/English) across all API responses
        and the admin panel.
  \item Implement offline mode with client-side caching for saved books and
        notes.
  \item Add direct Q\&A for course lesson transcripts using the existing RAG
        infrastructure.
\end{itemize}

% ===========================
\section*{Chapter Summary}
\addcontentsline{toc}{section}{Chapter Summary}

This chapter documented the backend implementation of Learnova --- a robust,
secure, and AI-integrated system built on Laravel 12, FilamentPHP 3, MySQL 8.0,
and a suite of 6 Python AI microservices. The proxy architecture ensures that
the platform degrades gracefully when AI services are unavailable, maintaining
core functionality at all times. The \texttt{TranscribeAllLessons} command
demonstrates the system's extensibility, enabling AI features to be added
incrementally to existing content modules. The admin panel provides
administrators with a powerful graphical interface for managing all platform
content. The next chapter details the intelligent AI features built on top of
this backend infrastructure.